\documentclass[11pt,a4paper]{article}
\usepackage[T1]{fontenc}
\usepackage{lmodern,amsmath,amssymb,amsthm,mathtools}
\usepackage[margin=26mm]{geometry}
\usepackage{array,longtable,booktabs,xcolor,microtype}
\usepackage[colorlinks=true,linkcolor=blue!45!black,citecolor=blue!45!black,urlcolor=blue!45!black]{hyperref}
\newtheorem{theorem}{Theorem}[section]
\newtheorem{lemma}[theorem]{Lemma}
\newtheorem{proposition}[theorem]{Proposition}
\newtheorem{corollary}[theorem]{Corollary}
\theoremstyle{definition}
\newtheorem{definition}[theorem]{Definition}
\theoremstyle{remark}
\newtheorem{remark}[theorem]{Remark}
\newcommand{\Q}{\mathbb Q}
\newcommand{\Z}{\mathbb Z}
\newcommand{\Sha}{\operatorname{Sha}}
\setlength{\emergencystretch}{2em}
\setcounter{tocdepth}{1}
\title{Small-variable polynomial constructions over the rationals}
\author{}
\date{7 November 2026}
\begin{document}
\maketitle
\begin{abstract}
We construct integer polynomials in 29 and 31 total variables
by eliminating finite algebraic choices from an explicit arithmetic
system. Their degrees are bounded by $2^{2093605}$ and $2^{180}$,
respectively. We prove the polynomial constructions, their numerical
bounds, the local input normalization and an exact index-depth formula.
Under four precisely stated arithmetic inputs their rational input
projections equal the integers. Every finite expression is specified
by reproducible data, with independent algebraic checks.
\end{abstract}
\tableofcontents
\newpage
\section{The constructions and their scope}
For a polynomial $P\in\mathbb Z[T,Y_1,\ldots,Y_m]$, its rational
input projection is the set of $t\in\mathbb Q$ for which some rational
tuple $\mathbf y$ satisfies $P(t,\mathbf y)=0$. We count $T$ among the
variables. The degree is the total degree, and the coefficient height
is the largest absolute value of an integer coefficient.

We give two complete finite specifications with 29 and 31 total
variables. Their polynomial character and numerical bounds are proved
by finite algebra. The assertion that their input projections equal
$\mathbb Z$ is a corollary of the four arithmetic inputs in
Section~\ref{sec:u29}. Those inputs are stated in full and attributed
to their sources. They are not inferred from the coefficient checks.

The construction separates retained rational coordinates from finite
algebraic choices. Monic equations for the latter define a finite free
algebra. A separating linear form and a reciprocal equation then
replace all its finite choices by one additional rational coordinate.
The resulting formulas are actual integer polynomials: their
determinants and finite products do not introduce quantified roots.

\section{The exact numerical height constant}
\label{sec:exact-height-constant}

The contact-height input uses a fixed, explicitly specified integer.
Its exact evaluation improves the numerical bounds in both construction
methods. Put
\[
 B_*=\prod_{p\le76176}p,\qquad
 H_{19}=2^{110}210^{1800}B_*^{19}+128.
\]
This is the constant of the source's exponent-nineteen height theorem
\cite[Theorem \nolinkurl{thm:quartic-height19}]{RepositoryV9}.

\begin{lemma}\label{lem:exact-height-constant}
The exact inequalities
\[
                 2^{2093447}<H_{19}<2^{2093448}
\]
hold. In particular one may use $h=2093448$ in every argument
requiring $H_{19}<2^h$.
\end{lemma}
\begin{proof}
This finite integer calculation has the following reproducible
certificate. Enumerate primes through $76176$ by marking, for
each unmarked $p\le\lfloor\sqrt{76176}\rfloor$, all multiples
starting at $p^2$. The remaining integers at least two are
exactly the primes: every composite has a prime divisor not
exceeding its square root, and no prime is marked.
An independent enumeration tests every integer by trial division
through its integer square root. Both lists have $7499$ entries
and agree entry by entry.

Their exact product has $109445$ binary digits. Computing the
displayed powers and product as integers gives $2093448$ binary
digits for $H_{19}$. The comparison with the next smaller power
of two is strict as well. The supplied checker performs both
prime enumerations and both integer comparisons; the complete
prime list and the unsigned binary encoding of $H_{19}$ are
included. No floating-point logarithm enters these comparisons.
\end{proof}

This lemma evaluates a constant. It does not prove the
contact-height inequality itself. The latter is still the
arithmetic input stated in Section~\ref{sec:u29}, now used with
this exact value and its sharper certified bound.


\section{An explicit polynomial for a guarded finite algebra}
\label{sec:explicit-finite-algebra}

Let $R=\mathbb Z[x_1,\ldots,x_b]$. Consider a triangular monic presentation
\[
 A=R[z_1,\ldots,z_m]/(f_1,\ldots,f_m),\qquad
 f_i\in R[z_1,\ldots,z_i],
\]
where $f_i$ is monic of degree $e_i\ge1$ in $z_i$. Its standard
monomials $\prod_i z_i^{a_i}$, $0\le a_i<e_i$, form an $R$-basis.
In particular $A$ is free of rank $n=\prod_i e_i$. For $a\in A$,
write $M_a$ for its multiplication matrix in this ordered monomial basis.
The matrices $M_a$ commute, and their entries are integer polynomials in
the base coordinates. They are completely specified by repeated reduction
using the monic equations; they need not be materialized to specify a
determinant involving them.

Fix elements $G,E_1,\ldots,E_r\in A$. We assume that in every geometric
fiber, every point at which $G$ is nonzero has local length one. This
condition holds, for example, if $G$ contains the product of the diagonal
derivatives $\partial f_i/\partial z_i$ as a factor. We seek the rational
base points admitting a rational point of $A$ with
\[
 G\ne0,\qquad E_1=\cdots=E_r=0.
\]

For $0\le\ell\le m n(n+1)/2$, define
\[
 \lambda_\ell=1+\sum_{i=1}^m\ell^i z_i,\qquad
 M=M_{G\lambda_\ell},\qquad
 H(V)=\det(I-VM)=\sum_{k=0}^n a_kV^k.
\]
Here and below $a_k\in R$ depend on the fixed integer $\ell$; always
$a_0=1$. Also put
\[
 R_i(V)=\operatorname{Tr}\bigl(M_{E_i}\operatorname{adj}(I-VM)\bigr).
\]
This is a polynomial in $R[V]$ of degree at most $n-1$.

For each $1\le d\le n$, set
\begin{align*}
 H_d(V)&=\sum_{k=0}^d a_kV^k,\\
 g&=a_d\operatorname{Res}_{V}(H_d,H'_d),\qquad L=a_d^2g,\\
 \Psi(U)&=1+\sum_{j=1}^d a_{d-j}a_d^{2j-1}g^jU^j,\\
 \widetilde R_i(U)&=(LU)^{n-1}R_i\bigl(1/(LU)\bigr),\\
 Q_{\ell,d}(\mathbf x,U)&=
 \Psi(U)^2+\sum_{k=d+1}^n a_k^2+\sum_{i=1}^r\widetilde R_i(U)^2.
\end{align*}
The resultant is the Sylvester determinant with the specified degree
bounds $d$ and $d-1$. The expression for $\widetilde R_i$ denotes
polynomial coefficient reversal, hence defines a polynomial also at
$LU=0$. Every displayed expression has integer coefficients.

\begin{theorem}\label{thm:explicit-guarded-finite-algebra}
The explicitly specified polynomial
\[
 P(\mathbf x,U)=
 \prod_{\ell=0}^{m n(n+1)/2}\ \prod_{d=1}^{n} Q_{\ell,d}(\mathbf x,U)
\]
has the following property. For every $\mathbf a\in\mathbb Q^b$,
there exists $U\in\mathbb Q$ with $P(\mathbf a,U)=0$ if and only if
there is a rational point of the fiber algebra $A_{\mathbf a}$ satisfying
$G\ne0$ and all $E_i=0$. The formula uses only one coordinate in addition
to the $b$ base coordinates; its formal algebra generators are eliminated
by the determinants and are not quantified coordinates.
\end{theorem}
\begin{proof}
Fix a rational base point. First consider one factor $Q_{\ell,d}$.
At $a_dg=0$ or $U=0$, its reciprocal residual $\Psi$ equals one.
Otherwise, with $V=1/(LU)$,
\[
 \Psi(U)=(LU)^d\frac{H_d(V)}{a_d}.
\]
Thus a rational zero of $Q_{\ell,d}$ implies $a_dgU\ne0$, all
coefficients $a_k$ with $k>d$ vanish, and $H(V)=H_d(V)=0$.
The resultant condition makes this a simple root of $H$. Since
$H(0)=1$, the root $V$ is nonzero.

The rational matrix $I-VM$ has a one-dimensional kernel. Every
multiplication matrix $M_{z_i}$ preserves it and therefore acts on
it by a rational scalar $\zeta_i$. Applying the algebra relations
to this common eigenline shows that $(\zeta_1,\ldots,\zeta_m)$ is
a rational point of the fiber. Its $G\lambda_\ell$-value is $1/V$,
so its $G$-value is nonzero.

To verify the closed conditions, write
$J=\operatorname{Tr}(\operatorname{adj}(I-VM))$. At a root of $H$,
the adjugate identity gives
$M\operatorname{adj}(I-VM)=V^{-1}\operatorname{adj}(I-VM)$.
Differentiating the determinant then yields
\[
 J=-V H'(V)\ne0.
\]
The columns of the adjugate lie in the common eigenline, so
\[
 R_i(V)=E_i(\boldsymbol\zeta)J.
\]
Since $LU\ne0$, the equation $\widetilde R_i(U)=0$ is equivalent
to $R_i(V)=0$, and hence $E_i(\boldsymbol\zeta)=0$. This proves
the reverse implication without assuming that the rest of the fiber
algebra is reduced.

Conversely, suppose a suitable rational point exists. Over an
algebraic closure, multiplication by $G\lambda_\ell$ has the
value $G(p)\lambda_\ell(p)$ at every geometric point $p$, with
multiplicity equal to its local length. If $G(p)=0$, its factor
in $\det(I-VM)$ is one, even if the local algebra is nonreduced.
If $G(p)\ne0$, its local length is one by hypothesis.

At any such point $p$, the polynomial
$G(p)(1+\sum_{i=1}^m\ell^i z_i(p))$ is nonzero as a polynomial
in $\ell$, because its constant coefficient is $G(p)\ne0$.
For distinct such points $p,q$, their difference is also nonzero:
if all coefficients vanished, first $G(p)=G(q)\ne0$, and then
all their coordinates would agree. Each of these polynomials has
degree at most $m$. With at most $n$ geometric points, there are
at most $m n(n+1)/2$ bad values of $\ell$ in total, counting zero
values and collisions. Thus one of the specified integers makes
all nonzero guarded eigenvalues distinct and nonzero.

For that $\ell$, the polynomial $H$ is squarefree of some degree
$1\le d\le n$. The given rational point provides a nonzero
rational root $V$ of $H$. Its closed residuals vanish, so the trace
identities give $R_i(V)=0$. The leading coefficient and resultant
are nonzero. Taking $U=1/(LV)$ therefore gives a zero of the
corresponding factor $Q_{\ell,d}$, and hence of $P$.
\end{proof}

The degree cases are part of the literal finite formula. They include
every exceptional drop in degree of $H$. The discriminant condition
excludes bad separating choices, and the finite bound on $\ell$ proves
that this exclusion loses no valid rational point. In particular the
construction does not leave a generic primitive-element argument or an
unperformed exceptional-stratum decomposition implicit.


\section{Odd products and a 29-variable construction}
\label{sec:u29}

Variable counts include the distinguished input.  The construction uses
Conant's odd-product height identity and direct-target comparison argument
as presented in \cite[Lemma \nolinkurl{v9:lem:product} and
Lemma \nolinkurl{lem:v9-bridge}]{RepositoryV9}, together with the explicit
finite algebra construction.  We give the rational system, prove its
input projection under precisely stated arithmetic hypotheses, and then
eliminate its finite algebraic choices.  The product identity and point
charts are proved independently of the contact-height and prime-pattern
inputs.

\subsection{Arithmetic inputs and conventions}

Put
\[
 K=\mathbb Q(\sqrt2),\quad d=3-\sqrt2,\quad D=d^2=11-6\sqrt2,
 \qquad E:y^2=x^3-Dx,
\]
and fix
\[
 R_0=(1+\sqrt2,2\sqrt2),\qquad
 P=2R_0=\left(\frac{57}{8}-\frac72\sqrt2,
          -\frac{165}{8}+\frac{417}{32}\sqrt2\right).
\]
Write $\sigma$ for conjugation, $T_0=(0,0)$, $f(Q)=1/x(Q)$,
and $z(Q)=-x(Q)/y(Q)$.  Heights are absolute logarithmic Weil
heights, with normalized local degree weights.  Valuations at
unramified primes satisfy $v(p)=1$.

We use the following arithmetic hypotheses, with the named source
locations in \cite{RepositoryV9}.
\begin{enumerate}
\item (Elliptic group and height.)\\
      Source: \nolinkurl{subsec:v5-point-source}, \nolinkurl{eq:o8-indexheight}.\par $2E(K)=\mathbb ZP$, and
      $h(x([n]P))\ge n^2-6$ for every nonzero integer $n$.
\item (Contact height: Theorem \nolinkurl{thm:quartic-height19}.)
      For $t=u/v$ in lowest terms, $v>0$, and
      $t\notin\mathcal D=\{0,1,2401,-343/32\}$, define $M_5(t)$
      to be the product of the primes $p>76176$ for which at least one of
\[
 v_p(v),\quad v_p(u),\quad v_p(u-v),\quad
 v_p(u-2401v),\quad v_p(32u+343v)
\]
      is positive and odd.  These are the normalized contacts with
      $\infty,0,1,2401,-343/32$.  Put
\[
 B_* =\prod_{p\le76176}p,\qquad
 H=H_{19}=2^{110}210^{1800}B_*^{19}+128.
\]
      The inherited assertion is $h(t)\le H M_5(t)^{19}$.
      Exact integer evaluation gives $H<2^{2093448}$.
      In either profile below we may use the indicated weaker exponent $\eta$.
\item (Pole obstruction and realization.)\\
      Source: \nolinkurl{hyp:o8-pole} and
      \nolinkurl{lem:original-seven-simultaneous-realization}.
      The pole equations, with $B=17b$,
\[
 U=b\frac{1-bt^2}{1+bt^2},\quad h_0=e(B-U^2),\quad
 p^2=B(1-e)(1+3e),\quad
 Bq^2=(s^2-Bh_0)(s^2+3Bh_0),
\]
      exclude negative odd $v_\ell(b)$ for $\ell\notin\{2,3,17\}$.
      They are soluble with $t=0$ whenever $b$ has square reduced
      denominator, is positive and less than $17$, is square in
      $\mathbb Q_2,\mathbb Q_3,\mathbb Q_{17}$, and satisfies
      $v_2(b)\ge3$, $v_3(b)\ge1$, $v_{17}(b)\ge2$.
\item (Simultaneous prime patterns: Theorem \nolinkurl{thm:v5-simultaneous-primes},
      Lemma \nolinkurl{lem:v5-two-reductions}, and the application of
      \cite[Theorem 4.1]{Sun2026Rational86}.)
      The simultaneous Gaussian prime-pattern input supplies the
      prescribed prime values $u+c_jw^2$, arbitrary fixed divisibility
      of $w$, and arbitrarily large $u/w^2$.  Their two conjugate
      reduction orders satisfy
\[
 v_2(d_-)=6,\quad v_2(d_+)=2,\quad
 \gcd(d_-(p),d_+(p'))=4
\]
      for every pair of selected primes.  An arbitrary fixed odd
      avoidance modulus can be imposed.  For each selected Gaussian
      prime $\pi$, the source convention is
      $\pi\equiv1+3i\pmod7$, making both $N(\pi-1)$ and
      $N(\pi+1)$ units at seven.  At each other odd prime $\ell$
      in the fixed avoidance modulus, the condition
      $\pi\equiv i\pmod\ell$ makes both orders congruent to two
      modulo $\ell$.  Thus the orders avoid the entire odd part
      of the fixed modulus.  The selected rational primes are $1$
      modulo $8\cdot3\cdot17$.
\end{enumerate}
These hypotheses include the simultaneous application of Sun's stated
prime-pattern theorem.  They are assumptions of the theorem below,
not conclusions of its implementation audit.
The semilocal index-size condition is a separate elementary consequence
of the displayed curve and point, proved in Proposition~\ref{anchor3:exact-depth};
it is not an additional arithmetic hypothesis.

The two final profiles are fixed as follows.  The symbols $\mu,\kappa$
are assigned divisibility orders, while $\rho_0,\rho_1$ are actual formal
logarithm truncation precisions.  The scalar ratio order is $m$.
\[
\begin{array}{c|c|cc|cc|c|c}
\text{variables}&\eta&\mu&\kappa&\rho_0&\rho_1&m&\text{size certificate}\\\hline
29&24&24&128&24&128&100485552&M=2^{2093449}\\
31&19&10&39&12&40&20&C=3^{2^{26}}.
\end{array}
\]
For the first profile the height exponent 24 is a consequence of the
exponent-19 input.  No new arithmetic hypothesis is introduced by the
large scalar order.  Section~\ref{sec:u29-height-cutoff} proves the
size and precision inequalities for both profiles.

\subsection{An elementary Lucas growth estimate}

The following estimate is independent of the arithmetic hypotheses.
It supplies the growth inequality used in the integer-extraction argument.

\begin{lemma}\label{lem:lucas-growth-elementary}
Let $a,t,b\in\mathbb Z$, with $|a|>7$ and $b\ne0$, satisfy
\[
 t^2-a(a-2)tb+(a-2)^2b^2=1.
\]
Then $\log|t|>|a|-3$. In particular,
$\log\max(|a|,|t|,|b|)>|a|-3$.
\end{lemma}
\begin{proof}
Put $A=|a|$, $c=(a-2)b$ and $c'=\operatorname{sgn}(a)c$.
Then
\[
                       t^2-Atc'+(c')^2=1.
\]
We have $|c'|\ge|a-2|\ge6$, so $t\ne0$. The two terms $t,c'$
have the same sign: opposite signs would make the left-hand side
at least $t^2+(c')^2>1$. Thus $T=|t|$ and $C=|c'|$ are positive
integer solutions of $T^2-ATC+C^2=1$.

Define $U_0=0$, $U_1=1$ and $U_{j+1}=AU_j-U_{j-1}$.
Every positive solution consists of adjacent terms of this sequence.
Indeed, order its coordinates as $M\ge m>0$. The integer
\[
                     Am-M=\frac{m^2-1}{M}
\]
lies in $[0,m)$, and replacing $(M,m)$ by $(m,Am-M)$ preserves
the equation. Repeated descent ends at $(1,0)$. Reversing the
steps gives the unordered pairs $\{U_{j+1},U_j\}$. Consequently
$C=U_j$ for some $j\ge1$ and $T\ge U_{j-1}$.

Put $d=|a-2|$. Since $d\mid C$, reduction of the recurrence
modulo $d$ gives $d\mid j$: if $a>0$, then $A\equiv2\pmod d$
and $U_j\equiv j\pmod d$; if $a<0$, then $A\equiv-2\pmod d$
and $U_j\equiv(-1)^{j-1}j\pmod d$. Hence $j\ge d\ge A-2$.
The recurrence also gives $U_r>(A-1)^{r-1}$ for $r\ge2$.
This follows by induction from $U_2=A>A-1$ and
$U_{r+1}>(A-1)U_r$, since $U_r>U_{r-1}$.
As $j\ge6$, we obtain
\[
 \log T>(j-2)\log(A-1)\ge(A-4)\log7>A-3
 \qquad(A\ge8).
\]
The last inequality holds at $A=8$ and its difference increases
with $A$. This proves the assertion.
\end{proof}


\subsection{An exact height identity}

\begin{lemma}\label{lem:u29-product}
For any nonzero odd integers $n_1,\ldots,n_r$,
\[
 h\left(\prod_{i=1}^r x([n_i]P)\right)
       =\sum_{i=1}^r h(x([n_i]P)).
\]
This identity requires only that the displayed $P=2R_0$ have infinite
order; it does not require the contact-height or prime-pattern hypotheses.
\end{lemma}
\begin{proof}
First consider a finite place $v$ with $v(D)=0$.  The formal subgroup
$E_1(K_v)$ consists of the identity and points with $v(x)<0$.
It is an open subgroup, so $E(K_v)/E_1(K_v)$ is finite.  These facts
hold for the displayed integral Weierstrass equation, including the
place above two.  The image of $T_0$ has order two, since
$T_0\notin E_1(K_v)$, and
\[
                    x(Q+T_0)=-D/x(Q).
\]
If an odd multiple of $P$ belongs to $E_1(K_v)$, the image of $P$
in the quotient has odd order.  No multiple of that image is the
nonzero element represented by $T_0$.  Hence no odd multiple of $P$
has positive $x$-valuation.  If no odd multiple belongs to $E_1(K_v)$,
all odd multiples have nonnegative $x$-valuation.  Thus the numbers
$v(x([n_i]P))$ have a common weak sign.

The only remaining finite place is the prime $(3-\sqrt2)$ above seven.
Here $v(D)=2$, and the discriminant has valuation six, so the displayed
integral model is minimal.  Its reduction is $y^2=x^3$.
The displayed $P$ reduces to $(1,6)$, a nonsingular point.  The subgroup
of points with nonsingular reduction therefore contains every multiple
of $P$.  A point with positive $x$-valuation would have positive
$y$-valuation and reduce to the singular point $(0,0)$.  Consequently
every multiple of $P$ has $x$-valuation at most zero at this place.

At either real embedding every $[n_i]P$ is a nonzero double.  The
duplication identity
\[
 x(2Q)=\left(\frac{x(Q)^2+D}{2y(Q)}\right)^2
\]
makes its abscissa positive.  The curve equation and its nonzero
ordinate imply $x([n_i]P)>d>1$ in each embedding.  Each real local
height therefore adds under multiplication.  At a finite place the
common weak sign similarly gives
\[
 \max\{0,-\textstyle\sum_i v(x([n_i]P))\}
       =\sum_i\max\{0,-v(x([n_i]P))\}.
\]
Summing with the normalized local weights proves the identity.
\end{proof}

\begin{corollary}\label{cor:u29-scales}
Let $Q_1,\ldots,Q_4$ be odd multiples of $P$ and write
\[
 C=\prod_{i=1}^4x(Q_i)=C_0+C_1\sqrt2.
\]
The three numbers $A_0=1$, $A_1=C_0^2$, and
$A_2=(1+C_1^2)^2$ are nonzero positive rational squares.
If the three contact detectors at $A_j\lambda$ imply
$M_0(A_j\lambda)^k\mid N\ne0$, then
\[
                         h(x(Q_i))<7H|N|^{\eta/k}
                         \qquad(1\le i\le4).
\]
Here $M_0$ is the odd-contact radical for
$\Theta=\{0,-1/2,-1,1/2399,-32/407\}$.
\end{corollary}
\begin{proof}
The product is positive in both embeddings, so $C_0>0$.
The fractional linear map $t\mapsto2+1/t$ identifies the five
contacts with the original five, outside the exceptional primes.
Consequently $h(t)\le(H+\log3)M_0(t)^{\eta}$.  Put
$T=(H+\log3)|N|^{\eta/k}$.  The unscaled detector and each scaled
detector give $h(A_j)\le2T$.  Hence $h(C_0)\le T$.
For $C_1=u/v$ in lowest terms with $v>0$,
$h(1+C_1^2)=\log(u^2+v^2)\ge2h(C_1)$, so $h(C_1)\le T/2$.
The elementary height inequality in the fixed basis gives
\[
 h(C)\le\tfrac32T+\tfrac32\log2<7H|N|^{\eta/k}.
\]
Lemma~\ref{lem:u29-product} bounds every summand by $h(C)$.
\end{proof}

\subsection{Point charts with no free parity witnesses}

\begin{lemma}\label{lem:u29-chart}
Suppose $R=(a,b)\in E(K)$ is not a two-torsion point.  At least one
of $R,R+T_0$ has nonzero rational component of its ordinate.
Thus every nonzero double is represented by $2R$ in the chart below,
and every odd multiple $[n]P$ with $n\ne1$ is represented by
$P+4R$ in the same chart.
\end{lemma}
\begin{proof}
If both ordinates had zero rational component, then
$b\in\sqrt2\mathbb Q^\times$ and
$D/a^2=b(R+T_0)/b(R)\in\mathbb Q$.  Therefore
$a/d\in\mathbb Q\cup\sqrt2\mathbb Q$.  Since
$b^2=d^3(a/d)((a/d)^2-1)\in\mathbb Q^\times$, this would force
$d^3\in\mathbb Q\cup\sqrt2\mathbb Q$, contrary to
$d^3=45-29\sqrt2$.  Translation by $T_0$ preserves both $2R$ and
$4R$.

If $n$ is odd and $n\ne1$, take $R=[(n-1)/2]R_0$, so
$P+4R=[n]P$.  Its $4R$ is nonzero.  The affine addition denominator
$x(4R)-x(P)$ cannot vanish: otherwise $4R=\pm P$, whereas
$4E(K)=2\mathbb ZP$ contains only even multiples.  The same
translation by $T_0$ ensures the chosen chart is defined.
\end{proof}

For each of the nine names $A,S,T,Z,H_0,U_A,U_S,U_T,U_Z$, retain a
rational pair $(a_0,a_1)$.  Write $a=a_0+a_1\sqrt2$ and
$a^3-Da=A'+B'\sqrt2$.  Introduce a finite algebra root $Y$ with
\[
 Y^4-4A'Y^2+8(B')^2=0,\qquad
 b=Y/2+(B'/Y)\sqrt2.
\]
Require $Y\ne0$, $(A')^2-2(B')^2\ne0$, and
$4Y^3-8A'Y\ne0$.  Then $R=(a,b)\in E(K)$.
For the four root names set the output point to $P+4R$; for the other
five set it to $2R$.  All arithmetic is by the usual affine addition
and duplication formulas with their denominators guarded.  The
derivative condition is automatic at rational points of this chart
but is retained to ensure geometric simplicity in the finite algebra.
Indeed, since $Y\ne0$, a vanishing derivative would give $Y^2=2A'$;
the quartic would then give $(A')^2=2(B')^2$, forcing $A'=B'=0$
over $\mathbb Q$, contrary to the norm guard.
It introduces no extra coordinate.

In particular, put $C_S=S-P=4R_S$.  This point, rather than $S$,
will be used in the scalar and size certificates.  Write
\[
 (A,S,T,Z)=([n]P,[a]P,[t]P,[b]P).
\]
The four labels are nonzero odd integers; the labels are consequences
of the group hypothesis, not additional quantified variables.

\subsection{The rational system}

Retain the input $X$, the eighteen point-chart coordinates, one
rational $\lambda$, two parameters for each of three comparison poles,
and two parameters for a scalar pole.  There are
\[
                            1+18+1+6+2=28
\]
retained rational coordinates.  Set
\[
                 \beta=216X^4+450459.
\]
The following elementary input map eliminates every scalar correction
choice.

\begin{lemma}\label{lem:u29-input-map}
For rational $X$, the number $\beta=216X^4+450459$ is integral
if and only if $X$ is integral. It is greater than $2401$.
For integral $X$, its contact
$r=\beta(\beta-1)(\beta-2401)(32\beta+343)(2\beta-1)$
is a square in $\mathbb Q_2,\mathbb Q_3,\mathbb Q_{17}$,
with $v_2(r)=v_3(r)=2$ and $v_{17}(r)\in\{0,2\}$.
\end{lemma}
\begin{proof}
This is proved, with its complete residue table and exceptional
valuation calculation, in Lemmas~\ref{sn:integral-input}
and~\ref{sn:local-squares}.
\end{proof}

Put $Q_0=P$, $(Q_1,Q_2,Q_3,Q_4)=(A,S,T,Z)$,
$U_0=H_0$, and $V_i=U_i-Q_i$ for $0\le i\le4$.
For $k\in\{\rho_0,\rho_1\}$ define
\[
 \ell_k(Q)=z(Q)\sum_{j=0}^{k/4-1}c_j(Df(Q)^2)^j,
 \quad c_0=1,\quad c_1=-2/5,\quad
 c_j=c_{j-1}\frac{4j-5}{4j+1}.
\]
Put $(k_1,k_2,k_3,k_4)=(\mu,\kappa,\kappa,\kappa)$ and
$(r_1,r_2,r_3,r_4)=(\rho_0,\rho_1,\rho_1,\rho_1)$.  Set
\[
 L_i=\frac{\ell_{r_i}(U_i)}{\ell_{\rho_1}(H_0)},\qquad
 L'_i=\sigma\left(\frac{\ell_{r_i}(V_i)}{\ell_{\rho_1}(V_0)}\right).
\]
The sixteen field-valued tests, with their orders, are
\[
 \begin{array}{c|c}
 f(U_0),\ldots,f(U_4),\sigma f(V_0),\ldots,\sigma f(V_4)&1\text{ each}\\
 L_i-L'_i&k_i\\
 L_1-\beta&\mu\\
 L_3^2-L_2(L_2-2)L_3L_4+(L_2-2)^2L_4^2-1&\kappa.
 \end{array}
\]
Denote these tests by $D_j$ and their orders by $m_j$.
Use precisely the three scales of Corollary~\ref{cor:u29-scales}.
For each such $A_r$ put
\[
 t_r=A_r\lambda,\quad \delta_r=t_r-1,\quad
 F_r=\prod_{\theta\in\Theta}(t_r-\theta),\quad
 R_r=\frac{F_r}{A_r^5\delta_r^6},
\]
\[
 W_r=\prod_{j=1}^{16}
 \left(1+R_r\left(\frac{D_j}{\delta_r^{m_j+1}R_r^{m_j}}\right)^2\right),
 \qquad b_r=\frac{R_r}{\operatorname{Tr}_{K/\mathbb Q}(W_r^{-1})^2}.
\]
All $F_r$, $\delta_r$, norms of $W_r$, and displayed denominators
are required to be nonzero.  In the coefficient DAG, if
$W_r=w_0+w_1\sqrt2$, the last expression is exactly
\[
 b_r=\frac{F_r(w_0^2-2w_1^2)^2}
              {4A_r^5\delta_r^6w_0^2}.
\]

For each nonzero pole input $b$ use two retained parameters $\tau,s$
and a fixed finite choice $\varepsilon\in\{0,1\}$.  Set
\[
 B=17b,\quad e=\frac{\tau^2+(4\varepsilon-1)B}{\tau^2+3B},
 \quad h_0=e(B-b^2),\quad
 q^2=\frac{(s^2-Bh_0)(s^2+3Bh_0)}{B}.
\]
Require $\tau^2+3B\ne0$.  The first original conic is recovered by
$p=-4B\tau(1-\varepsilon)/(\tau^2+3B)$.
Conversely every point of that conic is covered: if $e\ne1$ take
$\tau=p/(e-1)$ and $\varepsilon=0$; if $e=1$ take
$\varepsilon=1,\tau=0$.  This includes $e=-1/3,p=0$.
There is no nonzero condition on $\tau,p$, or $q$.

For the fourth, scalar pole put
\[
 r_\beta=\beta(\beta-1)(\beta-2401)(32\beta+343)(2\beta-1),\quad
 \delta_\beta=1+(17r_\beta)^2,\quad
 R_\beta=\frac{102^2r_\beta}{\delta_\beta^2}.
\]
Require $R_\beta\ne0$.  Write the rational components of
$f(H_0),f(C_S),z(C_S)/z(H_0)$ as $g_0,\ldots,g_5$, in that order.
For $N_R(u,v)=u^2+Rv^2$, set
\[
 w_0=N_R(1,g_0/R),\quad
 w_1=N_R(g_1/R,g_2/R),\quad w_2=g_3/R,
 \quad w_3=N_R(g_4/R^m,g_5/R^m),
\]
\[
              W=N_R\bigl(N_R(w_0,w_1),N_R(w_2,w_3)\bigr),
              \qquad R=R_\beta.
\]
Require $W\ne0$, and give $R_\beta W^2$ the same pole certificate.
There is no scalar correction parameter.  The only original finite
choices are the four conic bits, hence $2^4$ source branches.
For rational input $R_\beta>0$, the constant-left norm tree in fact
has $W\ge1$, so this nonzero condition is automatic; it is harmless
and convenient to retain the same denominator-guard convention.

For the finite-anchor profile require
\[
 \prod_{j=1}^{M}\operatorname{Norm}_{K/\mathbb Q}
          \bigl(x(C_S)-x([j]P)\bigr)\ne0,
                  \qquad M=2^{2093449}.
\]
Clear denominators factor by factor using the canonical primitive triple
$x([j]P)=(a_j+b_j\sqrt2)/c_j$, $c_j>0$.
If $x(C_S)=(N_0+N_1\sqrt2)/D_0$, the integer factor is
\[
                  (c_jN_0-a_jD_0)^2-2(c_jN_1-b_jD_0)^2.
\]
The multiples of the fixed point are computed by exact binary elliptic
multiplication.  This is a finite, fully indexed polynomial guard,
although its coefficients are not expanded.

\subsection{Local implications and soundness}

We include the local argument to make the effect of direct targets
explicit.  At a prime $p>76176$ where $v_p(R)$ is odd, consider
$W=\prod_j(1+R(D_j/R^{m_j})^2)$, $T=\operatorname{Tr}(W^{-1})$,
and a pole certificate for $R/T^2$.  Each factor has valuation
\[
 \min\{0,v_p(R)+2(v_{\mathfrak p}(D_j)-m_jv_p(R))\},
\]
since its two terms have valuations of different parity.  If
$\lambda=\max_{\mathfrak p\mid p}v_{\mathfrak p}(W)\le0$, then
$v_p(T)\ge-\lambda$.  The pole implication gives
\[
                 0<v_p(R/T^2)\le v_p(R)+2\lambda.
\]
Thus $v_p(R)>0$ and $\lambda>-v_p(R)/2$.  At a place attaining
the maximum, every factor valuation is at least $\lambda$, whence
$v_{\mathfrak p}(D_j)>(m_j-3/4)v_p(R)\ge m_j-3/4$ and therefore
$v_{\mathfrak p}(D_j)\ge m_j$.
For the normalized detectors, $\delta_r$ is a unit at each consumed
contact prime, so this argument applies to the original $D_j$.

Write $H_0=[h]P$ and $U_i=[u_i]P$.  The first ten tests place the
two families in the corresponding formal groups.  Formal logarithms
have exact ratios $u_i/h$ and $(u_i-n_i)/(h-1)$.
The displayed truncations have relative precision $k_i$: the omitted
series terms start with valuation $k_i v(z)$, and the factorial bound
on their denominators gives the same conclusion for all subsequent
terms at $p>76176$.  At least one of $h,h-1$ is a $p$-unit.
The associated ratio is integral; the cross-test of order $k_i$
makes both ratios integral.  Relative precision is therefore absolute
precision.  Multiplication by $h$ and $h-1$ and subtraction gives
\[
                  L_i\equiv n_i\pmod{p^{k_i}}.
\]
This argument also holds at an inert prime.  Integer polynomial
congruences then give the corresponding congruence for the direct
Lucas residual.

The scalar certificate gives
\[
                 M_5(\beta)^m\mid(a-1),\qquad
                 h(\beta)\le H|a-1|^{\eta/m}.
\]
Indeed negative valuations of $\beta=216X^4+450459$
are multiples of four outside the fixed primes two and three.
At every consumed finite odd contact, $R_\beta$ has positive odd
valuation $\nu$; its normalizer is a unit and its fixed multiplier
is a square unit.
For $N_R$ there is no cancellation between the two terms.  If a leaf
at depth $j$ contributes to the root, its valuation has the form
$2^jv(\text{leaf})+e\nu$, $0\le e<2^j$.
The pole implication forces the root valuation to exceed $-\nu/2$;
hence every normalized leaf has valuation greater than $-\nu$.
An original test assigned order $m$ consequently has valuation
at least $(m-1)\nu+1\ge m$.  The order-one tests put $H_0,C_S$
in the formal group; the order-$m$ ratio implies
$v_p((a-1)/h)\ge m$, and thus $p^m\mid a-1$.
The contact-height hypothesis proves the stated height bound.

Set $c=|a-1|$, $B=\max(|n|,|a|,|t|,|b|)$ and
$\Delta=t^2-a(a-2)tb+(a-2)^2b^2-1$.
If $\Delta\ne0$, the local implications and the three product-height
scales give
\[
                    B^2-6<7H(5B^4)^{\eta/\kappa}.
\]
The index size supplied by either profile contradicts this inequality,
as proved in Section~\ref{sec:u29-height-cutoff}.  Hence $\Delta=0$
and $\log B>|a|-3\ge c-4$.
If $\beta\ne n$, let $N$ be the absolute value of the numerator of
$\beta-n$ in lowest terms.  It satisfies
$N\le2e^{h(\beta)}B$, and each contact radical has its $\mu$th power
dividing $N$.  Therefore
\[
                   B^2-6<7H(2e^{h(\beta)}B)^{\eta/\mu}.
\]
The scalar estimate and the same profile inequalities again contradict
Lucas growth.  Thus $\beta=n\in\mathbb Z$, and
Lemma~\ref{lem:u29-input-map} implies $X\in\mathbb Z$.

\subsection{Simultaneous completeness}

Fix $X\in\mathbb Z$.  Lemma~\ref{lem:u29-input-map} gives an odd
integer $\beta>2401$, with $r_\beta$ positive and square at the three
preparation primes.  The integer $\delta_\beta=1+(17r_\beta)^2$
is a unit at two and three because $r_\beta$ is divisible by both,
and at seventeen because of its factor seventeen.  Consequently
$R_\beta$ has square reduced denominator, is square locally, and has
valuations at least $4,4,2$ at $2,3,17$.  Moreover
\[
 0<R_\beta=\frac{102^2r_\beta}{(1+(17r_\beta)^2)^2}
       <\frac{102^2}{17^4}<1.
\]
No finite squareclass adjustment is required.

Choose a real compact neighborhood for $H_0$ near, but away from, the
identity in both embeddings, with $f(H_0)$ small and $z(H_0)$ bounded
away from zero.  The positive orbit of every suitable finite-index
subgroup of $\mathbb ZP$ is dense in the identity component of its
real closure.  Thus one may choose a positive $h_*$ divisible by
$64$, and deep in every formal group over the finite support of
$R_\beta$, whose point lies in this neighborhood.  Choose a positive
integer $s$ divisible by six, still deeper at these places, sufficiently large for the selected size certificate, and with $[s]P$ sufficiently close to the real identity.
Set $a=s+1$, so $a$ is odd and $a\equiv1\pmod3$.
Take $S=[a]P$, and use $C_S=[s]P$ in the scalar certificate.
The normalized leaves of its norm tree are then small in the real
absolute value and sufficiently deep locally, uniformly for $H_0$ in
a smaller neighborhood.  Its root is a local unit and
$|W-1|<1/2$.

Odd Lucas labels are available explicitly.  Let
$U_0=0$, $U_1=1$, $U_{j+1}=aU_j-U_{j-1}$, and put
\[
              k=2(a-2),\qquad b=U_k/(a-2),\qquad t=U_{k+1}.
\]
Modulo $a-2$ one has $U_j\equiv j$, so $b$ is integral.
Modulo two the sequence has period $0,1,1$.  Since $a\equiv1\pmod3$,
$k\equiv1\pmod3$, so $b,t$ are odd.  They are positive and greater
than one.  The identity
$U_{k+1}^2-aU_kU_{k+1}+U_k^2=1$ gives $\Delta=0$.
The four root labels $\beta,a,t,b$ are therefore all odd and at least
three, and Lemma~\ref{lem:u29-chart} covers their points.

Temporarily set $H_0=[h]P$ and $U_i=[hn_i]P$.  Local continuity
provides a progression $h\equiv h_*\pmod L$, with $64\mid L$,
preserving the scalar units and fixed lower bounds for all comparison
tests at $2,3,17$.  Include the order of the finite component group
of the real closure in $L$.  Remove the finitely many comparison
pole loci from the real neighborhood and pass to a compact subset;
the tests are then uniformly bounded there.  These choices precede
the prime-pattern application.

Choose $D_0$ with $D_0^2\theta/A_r\in\mathbb Z$ for every contact
and scale, put $v=D_0^2w^2$, and choose $\lambda=u/v$.
The contact numerators $u-(\theta/A_r)v$ are a finite set of distinct
prime forms after repetitions are removed.  Apply the simultaneous
prime-pattern hypothesis beyond all fixed bad primes, with enough
divisibility of $w$ and a sufficiently large lower bound on $u/w^2$.
The curve has full rational two-torsion at every selected residue
field.  Because $P$ is a double, $d_\pm(p)/4$ annihilates its
reduction.  Define
\[
 U=\operatorname{lcm}_p(d_-(p)/4),\qquad
 V=\operatorname{lcm}_p(d_+(p)/4).
\]
The order and avoidance hypotheses give
$\gcd(U,V)=1$, $\gcd(U,L)=16$, and $\gcd(V,L)=1$.
Thus
\[
             h\equiv0\pmod U,\qquad h\equiv1\pmod V,
             \qquad h\equiv h_*\pmod L
\]
are compatible.  Their progression still has the required real
closure, so $H_0$ can lie in the chosen compact set.  At one place over
each contact prime the points $H_0,U_i$ are formal; at the conjugate
place $V_0,V_i$ are formal.  Both exact logarithmic ratios are $n_i$.
All sixteen tests therefore have their assigned orders, including the
direct Lucas polynomial since $\Delta=0$.

For completeness of each comparison pole, note the exact identity
\[
 R_r=\frac{\prod_{\theta\in\Theta}(u-(\theta/A_r)v)}
                     {v^5(A_r\lambda-1)^6}.
\]
The denominator is a rational square.  At the five contact primes
$v_p(R_r)=1$.  At either place over such a prime each inverse norm
factor has residue zero or one, and at the successful place every
residue is one.  Its reciprocal trace is thus a unit with residue
one or two.  Consequently $b_r$ has positive odd valuation only at
these primes and has square reduced denominator.  Their prescribed
congruences make $b_r$ square at $2,3,17$.

Uniformly on the compact real set,
$R_r\sim A_r^{-6}\lambda^{-1}$,
$W_r=1+O(\lambda^{-3})$, and
$b_r\sim(4A_r^6\lambda)^{-1}$.
Hence $0<b_r<17$ for a sufficiently large prescribed lower bound on
$\lambda$.  At $p\in\{2,3,17\}$ put $\nu=v_p(v)$ and
$\eta=v_p(A_r)$.  Sufficient divisibility of $w$ gives
$v_p(R_r)=\nu-6\eta$.  The nonconstant terms of the norm factors
have valuations at least $3\nu+(10m_j-8)\eta-2c_p$, where $c_p$
is the previously fixed lower-bound constant.  Their inverse trace
therefore has valuation $v_p(2)$, and
$v_p(b_r)=\nu-6\eta-2v_p(2)$.  These valuations can exceed every
required preparation threshold.  The pole hypothesis supplies the
three comparison certificates.

The scalar pole input retains square denominator and the specified
local squareclasses because $W$ is a unit at every prime in its fixed
support.  Also $0<R_\beta W^2<4<17$.  The same pole hypothesis
supplies its certificate with preparation zero.  All guards are met,
and the conic parametrization covers all these witnesses.  This proves
completeness for every integer input.

\subsection{The finite polynomial and its bounds}

The thirteen formal roots consist of nine point quartics and four
pole quadratics.  A rational equation $q^2=N/D$ with $D\ne0$ is
made monic by $Z=Dq$, so $Z^2-ND=0$.  Its zero-radicand branch is
$Z=0$ with $ND=0$; its other branch has $Z^2=ND$ and $Z\ne0$.
Together with the point derivative guards this makes every admitted
geometric point simple.  The rank is at most
\[
                              4^9\,2^4=2^{22}.
\]
Theorem~\ref{thm:explicit-guarded-finite-algebra} eliminates all thirteen
roots using one new rational variable $U$.  For completeness, its
literal recipe is as follows.  On each quadratic support branch,
put $m=13$ and let $n=4^9 2^k$, where $k$ is the number of nonzero
quadratic branches.  Let $G$ be the product of the source guards and $E_i$ the closed
zero-radicand conditions.  For every
$0\le\ell\le mn(n+1)/2$ and $1\le d\le n$, form
\[
 \lambda_\ell=1+\sum_{i=1}^m\ell^iz_i,\quad
 H(V)=\det(I-VM_{G\lambda_\ell})=\sum a_jV^j,
 \quad R_i(V)=\operatorname{Tr}(M_{E_i}\operatorname{adj}(I-VM_{G\lambda_\ell})).
\]
Put $H_d=\sum_{j=0}^d a_jV^j$,
$g=a_d\operatorname{Res}(H_d,H'_d)$, $L=a_d^2g$, and take the factor
\[
 \left(1+\sum_{j=1}^d a_{d-j}a_d^{2j-1}g^jU^j\right)^2
       +\sum_{j>d}a_j^2
       +\sum_i\left((LU)^{n-1}R_i(1/(LU))\right)^2.
\]
Coefficient reversal makes the last expressions polynomials at
$LU=0$.  Multiply these factors over the displayed ranges, the
$2^4$ quadratic masks, and all $2^4$ original finite choices.
This is a completely specified integer polynomial in the retained
coordinates and $U$, with no formal root left as a free variable.

For the semilocal profile retain two additional rational coordinates
$v,w$.  With $f(C_S)=f_0+f_1\sqrt2$, put
$F_0=f_0^2+f_0f_1-f_1^2$ and $C=3^{2^{26}}$, and introduce the
finite quadratic root
\[
                 u^2=3F_0^2-6v^2+2w^2-2C^4.
\]
Lemma~\ref{lem:semilocal-three-witnesses}, together with the binary-form valuation in Lemma~\ref{lem:four-semilocal-inputs}, identifies this with
$f_0/C,f_1/C\in\mathbb Z_{(2)}\cap\mathbb Z_{(3)}$.
Proposition~\ref{anchor3:exact-depth} supplies the size bound used in the profile
proof.  Its zero root is included by the same quadratic support split.
For this profile the finite recipe has $m=14$, five quadratic support
bits, and branch rank $n=4^9 2^k$ with $0\le k\le5$; the four
original conic bits are unchanged.





\begin{theorem}\label{thm:u29}\label{thm:u29-preferred}
Under the stated arithmetic hypotheses, the two specified polynomials
$P_{29}$ and $P_{31}$ have rational input projection exactly $\mathbb Z$.
Their declared total numbers of variables, including the input, and
certified bounds are
\[
\begin{array}{c|c|c|c}
&\text{finite algebra rank}&\text{degree bound}&\text{coefficient-height bound}\\\hline
P_{29}&2^{22}&2^{2093605}&2^{\,2^{6280501}}\\
P_{31}&2^{23}&2^{180}&2^{\,2^{184}}.
\end{array}
\]
The coefficient height is the largest absolute value of an integer
coefficient.  The displayed degrees and heights are upper bounds,
not claims of exact expanded values.
\end{theorem}
\begin{proof}
The local implications and profile size calculations prove soundness.
The simultaneous construction proves completeness.  Every remaining
formal root has a monic equation, and all rational denominators are
explicitly guarded.  The finite algebra theorem therefore preserves
the rational input projection and uses one final coordinate beyond
the retained 28 or 30.  The coefficient DAGs have respectively 6386
and 4367 nodes.  Section~\ref{subsec:finite-algebra-bounds} proves
the degree and coefficient-height bounds from the finite algebra
and the fixed-point exclusion.
\end{proof}

The finite-anchor polynomial is specified by an exact indexed product
with canonically computed fixed coefficients.  The semilocal polynomial
has no such product.  Both use finite determinant and coefficient
operations; the construction gives complete polynomial specifications
without requiring expanded monomial lists.


\section{Quadratic height bounds for the fixed elliptic multiples}
\label{sec:quadratic-fixed-height}

Put $K=\mathbb Q(\sqrt2)$ and
\[
 D=11-6\sqrt2,\qquad E:\ y^2=x^3-Dx,\qquad
 P=\left(\frac{57-28\sqrt2}{8},
          \frac{-660+417\sqrt2}{32}\right).
\]
The height $h$ below is the absolute logarithmic Weil height. If
$x=(a+b\sqrt2)/c$ with $a,b,c\in\mathbb Z$, $c>0$, and
$\gcd(a,b,c)=1$, put
\[
 h_0(x)=\left\lceil\log_2\max(1,|a|,|b|,c)\right\rceil.
\]

\begin{lemma}\label{lem:fixed-height-nontorsion}
The point $P$ belongs to $E(K)$ and has infinite order.
\end{lemma}
\begin{proof}
Substitution verifies the curve equation. At the unique place above two,
normalize $v$ by $v(2)=1$. The displayed coordinates have valuations
$v(x(P))=-3$ and $v(y(P))=-9/2$, so the formal parameter
$z=-x/y$ has $v(z(P))=3/2$. The integral formal group of this
Weierstrass equation is odd: for every integer $m$,
$[m](Z)=mZ+Z^3g_m(Z^2)$ with integral coefficients.
At a parameter of valuation at least $3/2$, multiplication by an odd
integer preserves its valuation, and multiplication by two raises it
by one. In each case the leading term has strictly smaller valuation
than the remaining terms. Thus
\[
                   v(z([n]P))=3/2+v_2(n)
                   \qquad(n\ne0).
\]
It is finite, and hence $[n]P\ne\mathcal O$ for every nonzero $n$.
In particular all its abscissas are finite and nonzero.
\end{proof}

\begin{lemma}\label{lem:fixed-height-recurrences}
For $n\ge1$, write $h_n=h(x([n]P))$. Then
\[
 \begin{aligned}
 h_{2n}&\le4h_n+\log392,\\
 h_{2n+1}&\le2h_n+2h_{n+1}+\log8036.
 \end{aligned}
\]
The second formula is used for $n\ge1$.
\end{lemma}
\begin{proof}
The duplication and differential-addition identities are
\[
 x(2Q)=\frac{(x(Q)^2+D)^2}{4x(Q)(x(Q)^2-D)},
 \qquad
 x(Q+R)x(Q-R)=
       \left(\frac{x(Q)x(R)+D}{x(Q)-x(R)}\right)^2.
\]
For the second identity, the usual addition formulas reduce it to
\[
 (u+v)^2(uv-D)^2-4uv(u^2-D)(v^2-D)
                    =(u-v)^2(uv+D)^2,
\]
which is a polynomial identity. Taking $Q=[n]P$ and
$R=[n+1]P$ makes $x(Q-R)=x(P)$. Its denominator is nonzero:
an equality $x([n]P)=x([n+1]P)$ would contradict the infinite
order of $P$.

For each place, bound a homogeneous polynomial by the largest
coordinate norm to its degree. At a finite place the coefficient
bound is the maximum coefficient norm; at an infinite place use
the sum of the coefficient norms. Applied to duplication, whose
homogeneous numerator and denominator have degree four, this gives
\[
             h(x(2Q))\le4h(x(Q))+2h(D)+\log8.
\]
Applied to differential addition, with bidegree $(2,2)$ in the two
abscissas and denominator $x(P)(x(Q)-x(R))^2$, it gives
\[
 h(x(Q+R))\le2h(x(Q))+2h(x(R))+2h(D)+h(x(P))+\log4.
\]
For example, the archimedean coefficient sums for the second
formula are bounded by
$4\max(1,|D|)^2\max(1,|x(P)|)$, and the corresponding finite-place
bound has no factor four. These local estimates sum to the stated
absolute-height estimates.

Both real conjugates of $D$ exceed one, and $N_{K/\mathbb Q}(D)=49$,
so $h(D)=\log7$. Both real conjugates of $x(P)$ exceed one, its
norm is $41^2/8^2$, and its denominator ideal has norm $8^2$.
Consequently $h(x(P))=\log41$. Substitution gives
$2h(D)+\log8=\log392$ and
$2h(D)+h(x(P))+\log4=\log8036$.
\end{proof}

\begin{theorem}\label{thm:fixed-height-quadratic}
For every integer $n\ge1$,
\[
 h(x([n]P))\le n^2\log23534-\log574,
 \qquad h_0(x([n]P))\le30n^2.
\]
These inequalities are unconditional for the displayed curve and point.
\end{theorem}
\begin{proof}
Write $A=\log23534$ and $B=\log574$. The integer identities
\[
 23534=41\cdot574,\qquad
 574^3=23534\cdot8036,\qquad574^3\ge392
\]
give $h_1=A-B$, $3B=A+\log8036$, and $3B\ge\log392$.
Induction using Lemma~\ref{lem:fixed-height-recurrences} proves
$h_n\le An^2-B$. For an even index,
\[
 h_{2n}\le4(An^2-B)+\log392\le A(2n)^2-B.
\]
For an odd index at least three,
\[
 \begin{split}
 h_{2n+1}&\le2(An^2-B)+2(A(n+1)^2-B)+\log8036\\
 &=A(2n+1)^2+A-4B+\log8036
  =A(2n+1)^2-B.
 \end{split}
\]

We next compare coefficient height with absolute height. Write
$x=(a+b\sqrt2)/c$ in the primitive form above, and let
$\mathfrak d=\{u\in\mathcal O_K:ux\in\mathcal O_K\}$ be its
denominator ideal. Since $\mathcal O_K=\mathbb Z[\sqrt2]$, the
integer $c$ is the least positive rational integer with
$cx\in\mathcal O_K$. At a prime above a rational prime $p$, put
$d_{\mathfrak p}=\max(0,-v_{\mathfrak p}(x))$.
The $p$-adic exponent of $N\mathfrak d$ is
$\sum_{\mathfrak p\mid p}f_{\mathfrak p}d_{\mathfrak p}$.
Multiplying by a ramification index $e_{\mathfrak p}$ makes this
at least $d_{\mathfrak p}$. Therefore
$(N\mathfrak d)x$ is integral and $c\mid N\mathfrak d$.
If $h_{\mathrm{fin}}(x)$ denotes the finite-place part of $h(x)$,
then $N\mathfrak d=\exp(2h_{\mathrm{fin}}(x))$.

Writing $x'$ for the real conjugate, one has
\[
 |a|=c\left|\frac{x+x'}2\right|\le c\max(|x|,|x'|),\qquad
 |b|=c\left|\frac{x-x'}{2\sqrt2}\right|
                         \le c\max(|x|,|x'|).
\]
The product $\max(1,|x|)\max(1,|x'|)$ equals
$\exp(2h_{\infty}(x))$. It follows that
\[
                  \max(1,|a|,|b|,c)\le\exp(2h(x)).
\]
Finally $23534<2^{15}$, so the first height bound gives
$\exp(2h_n)<2^{30n^2}$ and hence $h_0(x([n]P))\le30n^2$.
\end{proof}

\subsection{The indexed exclusion polynomial}

Let $\mathcal A$ be a finite free algebra over an integer polynomial
ring, with a fixed basis. For a multiplication matrix use the maximum
column sum of the coefficient $\ell^1$-norms. This norm is
submultiplicative. Suppose $b\ge0$ and the multiplication matrices of
$N_0,N_1,D_0\in\mathcal A$ have norm at most $2^b$.
Write $x([n]P)=(a_n+b_n\sqrt2)/c_n$ in primitive form and set
\[
 \mathcal E_n=(c_nN_0-a_nD_0)^2-2(c_nN_1-b_nD_0)^2,
 \qquad \mathcal E_{\le M}=\prod_{n=1}^M\mathcal E_n.
\]
Each fixed triple is computed by exact elliptic arithmetic. The product
is finite and is a literal polynomial specification, independently of
whether its coefficients are expanded.

\begin{proposition}\label{prop:fixed-height-cubic-product}
For every integer $M\ge1$,
\[
 \|M_{\mathcal E_{\le M}}\|
       \le 2^{20M^3+30M^2+(2b+14)M}
       \le 2^{(2b+64)M^3}.
\]
If $b\le1048$ and $M\ge64$, the first exponent is at most $21M^3$.
\end{proposition}
\begin{proof}
Theorem~\ref{thm:fixed-height-quadratic} and the triangle inequality give
\[
 \|M_{c_nN_0-a_nD_0}\|,\ \|M_{c_nN_1-b_nD_0}\|
       \le2^{30n^2+b+1}.
\]
Squaring and summing the two terms of $\mathcal E_n$ therefore
gives $\|M_{\mathcal E_n}\|\le2^{60n^2+2b+4}$.
Sum these logarithmic bounds and use
$\sum_{n=1}^M n^2=M(M+1)(2M+1)/6$.
The second bound follows from $M,M^2\le M^3$.
For the last assertion, $2b+14\le2110$, and
\[
                   30\cdot64+2110=4030<64^2.
\]
The function $M^2-30M$ is increasing for integers $M\ge64$.
Thus $30M^2+(2b+14)M\le M^3$, proving the result.
\end{proof}

At a rational specialization with $D_0\ne0$, the factor
$\mathcal E_n$ vanishes exactly when
$(N_0+N_1\sqrt2)/D_0=x([n]P)$. Indeed a rational pair satisfies
$u^2-2v^2=0$ only when both entries vanish. Thus changing this
coefficient bound changes neither the exclusion polynomial nor its
rational zeros.

For the 29-variable specification, the recorded coordinate bound is
$b=1048$, and $M=2^{2093449}$. The indexed multiplication matrix
therefore has logarithmic norm at most
\[
                         21M^3<2^{6280352}.
\]
The independently checked determinant, reciprocal, trace and product
bounds consequently give
\[
            \deg P_{29}\le2^{2093605},\qquad
            H(P_{29})\le2^{\,2^{6280501}}.
\]
The source polynomial is specified by its finite-algebra relations,
branches, guards, and the indexed exclusion product above. Its
29 variables include the distinguished input. The 31-variable specification
does not contain this exclusion product and retains its bounds
$\deg P_{31}\le2^{180}$ and $H(P_{31})\le2^{\,2^{184}}$.
\section{Trading divisibility order for a shorter size exclusion}
\label{sec:u29-height-cutoff}

The number of factors in the finite size exclusion has a much larger
effect on the final degree than the powers in the scalar norm tree.
The two profiles use the explicit arithmetic input
\cite[Theorem \texttt{thm:quartic-height19}]{RepositoryV9}:
\[
 h(t)\le H M_5(t)^{19},\qquad
 1\le H<2^h,\qquad h=2093448.
\]
Here $H$ is the explicit integer
$2^{110}210^{1800}(\prod_{p\le76176}p)^{19}+128$, not an
unspecified constant.  Trial division enumerates exactly $7499$
primes in that product, the largest being $76163$.  Their product
has binary length $109445$, and exact integer substitution gives
binary length $2093448$ for $H$.  The supplied integer checker
recomputes these values without a supplied prime list or sieve.
The enumeration is complete because each candidate is divided by
every previously found prime not exceeding its square root; the
usual least-prime-divisor argument proves this criterion by induction.
The definition of binary length then gives the displayed strict bound.
It implies the same inequality with exponent 24.  We use that weaker
exponent for the finite-anchor construction below, and exponent 19
directly for the smaller-degree semilocal alternative.  No assertion
about additional prime separation is needed.

\subsection{A general parameter calculation}

Let the contact-height exponent be $\eta$.  Assign order $\kappa$
to the three Lucas cross-tests and to the Lucas polynomial, order $\mu$
to the input cross-test and input comparison, and order $m$ to the two
scalar ratio components.  The ten formal-membership tests retain order
one.  Choose formal logarithm truncations of relative precision at least
the assigned comparison orders.  The local, conic and norm-tree
arguments of Section~\ref{sec:u29} are unchanged.  If
\[
 c=|a-1|,\qquad B=\max(|n|,|a|,|t|,|b|),
\]
they give
\begin{align}
 h(\beta)&\le Hc^{\eta/m},\label{eq:u29-general-scalar}\\
 \Delta\ne0\quad&\Longrightarrow\quad
 B^2-6<7H(5B^4)^{\eta/\kappa},\label{eq:u29-general-Lucas}\\
 \beta\ne n\quad&\Longrightarrow\quad
 B^2-6<7H(2e^{h(\beta)}B)^{\eta/\mu}.
 \label{eq:u29-general-input}
\end{align}
For the second implication, $\Delta$ is the integral Lucas residual;
for the third, use the numerator of $\beta-n$ in lowest terms.
The third implication is needed only after the second has forced
$\Delta=0$, at which point
\[
                         \log B>|a|-3\ge c-4.
\]

An increase of $m$ does not require a longer formal logarithm series.
The scalar test is the valuation of $z(C_S)/z(H_0)$ itself.  For formal
points at the consumed primes, the logarithm divided by the local
parameter is a unit.  Thus
\[
 v_{\mathfrak p}\bigl(z(C_S)/z(H_0)\bigr)
       =v_p((a-1)/h_0),
\]
where $H_0=[h_0]P$.  This yields $p^m\mid a-1$ directly.  Larger $m$
only asks for deeper choices at finitely many places in completeness.
Since $m$ is fixed, the local continuity and Chinese remainder argument
still applies.  Powers such as $R_\beta^m$ are exact integer-power
operations in the coefficient DAG.

\begin{lemma}[A cutoff linear in the height constant]
\label{lem:u29-linear-cutoff}
Suppose $h\ge20$ is an integer and $1\le H<2^h$.  Work with the
contact-height exponent $\eta=24$, and set
\[
 E=h+1,\quad M=2^E,\quad m=48E,\quad
 \kappa=128,\quad\mu=24.
\]
If $c>M$, the three implications
\eqref{eq:u29-general-scalar}--\eqref{eq:u29-general-input}
force $\Delta=0$ and $\beta=n$.
\end{lemma}
\begin{proof}
Since $c$ and $a$ are integers, $|a|\ge c-1\ge M$.
Because $24/m=1/(2E)$ and the exponent $1/(2E)-1$ is negative,
\[
 \frac{Hc^{24/m}}{c}
 <2^h M^{1/(2E)-1}=2^{-1/2}<\frac34.
\]
Consequently $h(\beta)<3c/4$.

If $\Delta\ne0$, then $B\ge M>\sqrt{12}$ and
\[
 \tfrac12 B^2<B^2-6
 <7\,5^{3/16}HB^{3/4}<14HB^{3/4}.
\]
Thus $B^{5/4}<28H$.  On the other hand
\[
 B^5\ge2^{5h+5}>28^4\,2^{4h}>(28H)^4,
\]
since $h\ge20$ and $28^4<2^{20}$.  This is impossible, and
$\Delta=0$.

If $\beta\ne n$, equation~\eqref{eq:u29-general-input} becomes
$B^2-6<14He^{h(\beta)}B$.  Since $B,H,e^{h(\beta)}\ge1$,
\[
                 B<20He^{h(\beta)},\qquad
                 \log B<\log(32H)+h(\beta).
\]
For $h\ge5$ one has $(h+5)\log2<2^{h-2}$, as follows at
$h=5$ by direct comparison and thereafter by induction.  Therefore
\[
 \log(32H)<2^{h-2}=M/8<c/8,
 \qquad \log B<7c/8<c-4,
\]
the last inequality using $c>M\ge2^{21}$.
This contradicts Lucas growth.  Hence $\beta=n$.
\end{proof}

For $h=2093448$ the parameters are
\[
 M=2^{2093449},\qquad m=100485552,\qquad
 \kappa=128,\quad\mu=24.
\]
The longest formal logarithm truncation is 128, so the inherited
$p>76176$ estimate applies with ample room.  The first comparison
uses truncation 24.  Both directions of the rational source proof
are therefore valid with these parameters.  In particular the large
scalar exponent introduces no extra free coordinate, finite generator,
finite branch, or arithmetic hypothesis.

\subsection{Small test orders for the semilocal variant}

For a fixed 31-variable semilocal presentation, the literal anchor
constant has degree zero.  One can therefore reduce the degree by
decreasing the test powers and increasing that fixed constant.
Use the contact-height exponent 19 directly and assign
\[
                 \kappa=39,\qquad\mu=10,\qquad m=20.
\]
Use formal truncation 40 for the Lucas comparisons and denominators,
and truncation 12 for the input comparison.  A truncation precision
need only exceed the assigned divisibility order; they need not be
equal.  The inequality $\kappa>2\cdot19$ is what makes
\eqref{eq:u29-general-Lucas} subquadratic in $B$, and
$\mu>19/2$ is what makes \eqref{eq:u29-general-input} useful.
These are the smallest integer orders satisfying those two strict
inequalities.  This is a statement about this parameter calculation,
not an optimality theorem for all Diophantine definitions.

Choose the semilocal constant $C=3^{2^{26}}$.  Proposition~\ref{anchor3:exact-depth}
gives
\[
 c\ge2\,3^{2^{25}-1}>(64H)^{20}.
\]
For the displayed $h$, this follows from $3>2^{3/2}$ and the exact
integer comparison
\[
 1+\frac32(2^{25}-1)>20(h+6).
\]
Hence $h(\beta)\le Hc^{19/20}<c/64$.

If $\Delta\ne0$, the inequality $5^{19/39}<\sqrt5<16/7$ gives
\[
 B^2-6<16HB^{76/39}.
\]
As $B\ge|a|\ge c-1>\sqrt{12}$, it follows that
$B^{2/39}<32H$, and thus $B<(32H)^{39/2}$.
But $|a|\ge c-1>c/2>(64H)^{20}/2>(32H)^{39/2}$,
a contradiction.  We again obtain $\Delta=0$.

If $\beta\ne n$, then
\[
 B^2-6<7\,2^{19/10}He^{19h(\beta)/10}B^{19/10}.
\]
Dividing by $B^{19/10}$, and using $2^{19/10}<4$ and $B\ge1$,
gives $B^{1/10}<64He^{19h(\beta)/10}$ and therefore
\[
 \log B<10\log(64H)+19h(\beta)<\frac{29}{64}c<c-4.
\]
Here $\log(64H)<c/64$ follows, for example, from
$c>(64H)^{20}$ and $H\ge1$.
This contradicts Lucas growth, proving $\beta=n$.
Completeness only changes the finite local depths; the real estimate
for each normalized detector remains $1+O(\lambda^{-3})$ for every
positive test order.  Thus the same simultaneous construction applies.

\subsection{Effect on the explicit bounds}

These are the two profiles of Theorem~\ref{thm:u29-preferred}; they
alter no retained coordinate or finite algebra degree.  The
correction-free input map in Section~\ref{sec:u29} leaves only four
finite conic bits.  This reduces the finite product's original branch
count to $16$.  The coefficient checker incorporates that exact factor
count and the chosen size certificate.  The resulting bounds are
those in Theorem~\ref{thm:u29-preferred}.

The finite-anchor source is an exact indexed finite product.  Its
size exclusion has $2^{2093449}$ factors.  The coefficient-height
exponent has the form $3\log_2 M$ plus the recorded finite-algebra
overhead, which explains the role of shortening this exclusion.
The semilocal source instead uses one fixed power of three and has no
indexed exclusion.  The two profiles therefore serve different degree
objectives at their respective fixed numbers of variables.
\subsection{Explicit degree and coefficient bounds for the specification}
\label{subsec:finite-algebra-bounds}

The following bounds concern the expanded polynomial defined by the
indexed determinant formula. They do not require expanding its matrices
or its final finite product. For a polynomial $f$, let $\|f\|_1$ denote
the sum of the absolute values of its coefficients. Then
$\|fg\|_1\le\|f\|_1\|g\|_1$ and
$\|f+g\|_1\le\|f\|_1+\|g\|_1$.

Give every base coordinate weight one, and assign a positive integer
weight $w_i$ to each algebra generator $z_i$. Choose the weights so that
every nonleading term of its monic relation has weight at most
$e_iw_i$. Reduction by these relations never increases weighted degree.
Writing $W=\sum_i(e_i-1)w_i$, the entries of a multiplication matrix
$M_f$ therefore have base degree at most $\operatorname{wdeg}(f)+W$:
the input basis monomial has weight at most $W$, and the output basis
monomial has nonnegative weight.

For a matrix over the polynomial ring, use the maximum column sum of
the coefficient norms of its entries. This norm is submultiplicative.
If
\[
 z_i^{e_i}+\sum_{k<e_i}c_{ik}z_i^k=0,
\]
the block companion description gives
\[
 \|M_{z_i}\|\le
 \max\left\{1,\sum_{k<e_i}\|M_{c_{ik}}\|\right\}.
\]
Extending to later monic generators preserves this estimate, since
multiplication by an earlier element is repeated on diagonal blocks,
up to a simultaneous permutation of rows and columns. Thus weights
and bounds $\|M_{z_i}\|\le2^{b_i}$ can be computed recursively from
the coefficient DAG. The source compiler records the entire table,
and the independent checker verifies the corresponding inequalities.

For the point quartics the weights can be taken to be two. A coefficient
bound for the companion matrix is the sum of the bounds for its two
nonzero lower coefficients. For each of the subsequent quadratic
generators, its nonzero branch has only the constant coefficient
$LM$, so one may take
\[
 w_i\ge\tfrac12\operatorname{wdeg}(LM),\qquad
 b_i\ge\log_2\|M_{LM}\|.
\]
The zero branch has relation $z_i=0$ and satisfies the same bounds.
Consequently one uniform table applies to all $2^q$ fixed branches.

Put $n=2^{22}$ and $m=13$ for the finite-anchor construction, or
$n=2^{23}$ and $m=14$ for its semilocal variant, and
$T=m n(n+1)/2$ for the largest separating parameter. Suppose the
recorded bounds give
\[
 \operatorname{wdeg}(G)\le w_G,\quad \|M_G\|\le2^{b_G},\qquad
 \operatorname{wdeg}(E_i)\le w_E,\quad \|M_{E_i}\|\le2^{b_E}.
\]
Set
\begin{align*}
 b_\lambda&=\lceil\log_2(m+1)\rceil
       +\max_i\{i\lceil\log_2 T\rceil+b_i\},\\
 \alpha&=w_G+\max_i w_i+W,\qquad \beta=b_G+b_\lambda,\\
 \alpha_E&=w_E+W.
\end{align*}
These bound the entry degrees and operator norm of
$M_{G\lambda_\ell}$, uniformly for all prescribed $\ell$ and all
source branches.

The coefficients of $H(V)=\det(I-VM)$ have base degree at most
$A_H=n\alpha$, and its full coefficient norm is at most
$2^{B_H}$ with $B_H=n(\beta+1)$. Indeed the coefficient norm of a
determinant is bounded by the product of its column sums. This bound
already includes the sum over determinant permutations; no additional
factorial is needed. By multilinearity in columns,
\[
 R_i(V)=\left.\frac{\partial}{\partial t}
       \det(I-VM+tM_{E_i})\right|_{t=0}
\]
has base coefficient degree at most
$A_R=(n-1)\alpha+\alpha_E$ and coefficient norm at most $2^{B_R}$,
where
\[
 B_R=\lceil\log_2 n\rceil+b_E+(n-1)(\beta+1).
\]
This accounts for the $n$ possible columns containing $M_{E_i}$.

For each degree case $d\le n$, the Sylvester matrix for
$\operatorname{Res}(H_d,H'_d)$ has $2d-1$ rows. Every row has
coefficient norm at most $d\|H\|_1$. Hence uniform bounds for the
reciprocal scale $L=a_d^3\operatorname{Res}(H_d,H'_d)$ are
\begin{align*}
 A_L&=(2n+2)A_H,\\
 B_L&=(2n+2)B_H+(2n-1)\lceil\log_2 n\rceil.
\end{align*}
The reciprocal residual has degree at most $n(A_L+1)$ and coefficient
norm at most $2^{nB_L+\lceil\log_2(n+1)\rceil}$. The reversed trace
residuals have degree at most $A_R+(n-1)(A_L+1)$ and coefficient norm
at most $2^{B_R+(n-1)B_L}$. Together with the high-coefficient
conditions $a_k=0$, these give a bound for the sum of squares by
doubling the largest degree, doubling the largest logarithmic norm
bound, and adding $\lceil\log_2(n+q+1)\rceil$ to the latter.
For the reciprocal estimate, no division of coefficient norms is used:
its $j$th coefficient is the polynomial
$a_{d-j}a_d^{3j-1}\operatorname{Res}(H_d,H'_d)^j$.
Its degree is at most $jA_L$ and its logarithmic norm bound is at most
$j(3B_H+B_{\rm Res})=jB_L$.

The number of final factors is also explicit.  Put $n_0=4^9$ and
$q=4$ for the finite-anchor source, or $q=5$ for the semilocal
source.  For a mask with $k$ nonzero quadratic branches, the rank is
$n_k=n_0 2^k$.  There are $\binom qk$ such masks, each with
$n_k(1+mn_k(n_k+1)/2)$ separating-parameter and degree choices.
Including the $2^4$ original finite parameter choices gives exactly
\[
 B_{\mathrm{tot}}=2^4\left[
 n_0 3^q+\frac m2\bigl(n_0^3 9^q+n_0^2 5^q\bigr)\right]
\]
factors.  Thus the factor counts are respectively
\[
 12292021203320302469120,\qquad
 119138032421494864216064.
\]
Multiplying a uniform branch-degree bound and a uniform logarithmic
coefficient-norm bound by this factor count bounds the final product.

For the finite-anchor source the indexed guard is handled before these
operations.  Its coordinate weight is at most $129$, and its coordinate
operator logarithmic norm is at most $1048$.  With $M=2^{2093449}$,
Proposition~\ref{prop:fixed-height-cubic-product} gives
\[
 \operatorname{wdeg}(\mathcal E_{\le M})\le 2M\cdot129,
 \qquad
 \|M_{\mathcal E_{\le M}}\|\le 2^{21M^3}.
\]
All these are integer majorants.  The emitted power-of-two envelopes
for the guard, reciprocal scale, reversed trace and final factor are
checked successively by exact integer comparisons.  Their final values
are
\[
 \begin{array}{c|cc}
 &\deg P\text{ upper bound}&H(P)\text{ upper bound}\\ \hline
 P_{29}&2^{2093605}&2^{\,2^{6280501}}\\
 P_{31}&2^{180}&2^{\,2^{184}}.
 \end{array}
\]
The semilocal source records stronger unrounded integer majorants in
addition to these convenient envelopes.  These estimates concern the
expanded polynomial defined by the literal indexed recipe; no expanded
matrix or coefficient table is claimed to have been computed.
\section{Simultaneous parity and divisibility conditions}
\label{sec:parity-batching}

Put $K=\mathbb Q(\rho)$, where $\rho^2=2$, and consider
\[
 d=3-\rho,\qquad D=d^2=11-6\rho,\qquad
 E:\ y^2=x^3-Dx,
 \qquad
 P=\left(\frac{57-28\rho}{8},
          \frac{-660+417\rho}{32}\right).
\]
We write $v_p$ for the rational valuation with $v_p(p)=1$.
The valuation $\nu_2$ on $K$ above $2$ is normalized by
$\nu_2(\rho)=1$, so $\nu_2(2)=2$.
Let $R=\mathbb Z_{(2)}\cap\mathbb Z_{(3)}$, and set
\[
                         N(a,b)=a^2+ab-b^2.
\]

\begin{lemma}\label{lem:four-semilocal-inputs}
For rational $a,b,c,e$,
\[
 N(a,b)+6N(c,e)\in R
 \quad\Longleftrightarrow\quad a,b,c,e\in R.
\]
\end{lemma}
\begin{proof}
The binary form $N$ has no nontrivial zero over either $\mathbb F_2$
or $\mathbb F_3$. Consequently, for $p=2,3$ and $(a,b)\ne(0,0)$,
\[
                    v_p(N(a,b))=2\min(v_p(a),v_p(b)).
\]
Indeed, divide both inputs by the smaller power of $p$ and reduce
modulo $p$. The two summands $N(a,b)$ and $6N(c,e)$ have valuations
of opposite parity when both are nonzero. They therefore cannot
cancel. Their sum is $p$-integral if and only if
\[
 2\min(v_p(a),v_p(b))\ge0,
 \qquad 1+2\min(v_p(c),v_p(e))\ge0.
\]
Since these minima are integers, these inequalities say that all four
inputs are $p$-integral. If either pair is zero, the same conclusion
follows by omitting the corresponding zero term.
\end{proof}

\begin{lemma}\label{lem:semilocal-three-witnesses}
For $t\in\mathbb Q$,
\[
 t\in R\quad\Longleftrightarrow\quad
 \exists u,v,w\in\mathbb Q:\
                u^2+6v^2-2w^2-3t^2+2=0.
\]
\end{lemma}
\begin{proof}
This is the three-witness semilocal criterion used by the source
construction; we recall its arithmetic justification.
The Hilbert-symbol formulas show that $(2,3)_{\mathbb Q}$ is
ramified precisely at $2$ and $3$: outside $2,3$ both parameters
are units and the algebra splits; at $3$ the symbol is the
Legendre symbol $(2/3)=-1$, and reciprocity gives the symbol
$-1$ at $2$ since the real place is split. Its quadratic-field
embedding criterion \cite[Proposition 14.6.7]{Voight2021} says that a
nonzero rational number $\Delta$ is the square of a trace-zero
element precisely when $\Delta$ is nonsquare in both completions.
For a rational square $\Delta\ne0$, such an element is impossible
in a division algebra: $(q-\sqrt\Delta)(q+\sqrt\Delta)=0$ would
make $q$ scalar. For nonsquare $\Delta$, this is the usual local
embedding criterion for quadratic fields in quaternion algebras.

For $t\ne0$, take $\Delta=1-2/(3t^2)$, which is nonzero because
$2/3$ is not a rational square. The trace-zero square form is
$2x^2+3y^2-6z^2$. At $2$, if $v_2(t)>0$, then $v_2(\Delta)$ is
odd; if $v_2(t)=0$, then $\Delta\equiv3\pmod8$; and if
$v_2(t)<0$, then $\Delta\in1+8\mathbb Z_2$ is a square.
At $3$, a nonnegative valuation of $t$ gives odd negative valuation
of $\Delta$, whereas a negative valuation gives a square in
$1+3\mathbb Z_3$. Hence the embedding criterion holds exactly
when $t\in R$. Multiplying
$2x^2+3y^2-6z^2=\Delta$ by $3t^2$, and putting
$v=tx$, $u=3ty$, $w=3tz$, gives the displayed equation.
The case $t=0$ has witnesses $(u,v,w)=(0,0,1)$.
\end{proof}

\begin{lemma}\label{lem:exact-dyadic-index}
For every nonzero integer $n$,
\[
 \nu_2(x([n]P))=-6-4v_2(n).
\]
At the place above $3$, every odd multiple $[n]P$ has unit abscissa.
\end{lemma}
\begin{proof}
The displayed coordinates give
$\nu_2(x(P))=-6$, $\nu_2(y(P))=-9$, and therefore
$\nu_2(z(P))=3$ for the formal parameter $z=-x/y$.
The equation is integral at $2$. Its formal logarithm is an isometry
on the subgroup of depth strictly greater than
$\nu_2(2)/(2-1)=2$. No minimal-model assertion is needed for this
integral formal group \cite[Theorem IV.6.4]{Silverman2009}.
Its nonzero logarithm also proves that $P$ has infinite order.
Since the logarithm is a homomorphism,
\[
 \nu_2(z([n]P))=3+2v_2(n).
\]
The Laurent expansion $x=z^{-2}(1+O(z))$ then proves the formula.

At $3$ the equation has good reduction: its discriminant is
$64D^3$, a unit. The residue field is $\mathbb F_9$, and direct
reduction gives $\overline P=(\overline\rho,0)$.
This point has order two and nonzero abscissa. Every odd multiple
therefore has the same reduction and unit abscissa.
\end{proof}

\begin{theorem}\label{thm:parity-product}
Let $n_1,\ldots,n_r$ be nonzero integers and write
\[
        C=\prod_{i=1}^{r}x([n_i]P)=C_0+C_1\rho.
\]
Then
\[
 n_1,\ldots,n_r\text{ are all odd}
 \quad\Longleftrightarrow\quad
                2^{3r}C_0,\ 2^{3r}C_1\in R.
\]
\end{theorem}
\begin{proof}
The preceding lemma gives
\[
              \nu_2(2^{3r}C)=-4\sum_{i=1}^{r}v_2(n_i).
\]
Each summand on the right is nonnegative. Thus this expression is
nonnegative exactly when every index is odd.
The local integer rings at $2$ and $3$ are respectively
$\mathbb Z_2[\rho]$ and $\mathbb Z_3[\rho]$.
At $2$ this can also be seen directly from
\[
 \nu_2(a+b\rho)=\min(2v_2(a),2v_2(b)+1),
\]
whose two entries have opposite parity. At $3$, the integral basis
$(1,\rho)$ remains linearly independent modulo $3$.
Therefore integrality in either local ring is equivalent to
integrality of both rational coefficients. If the indices are odd,
the preceding lemma makes $C$ a $3$-adic unit, so the additional
$3$-integrality conditions impose no restriction. The converse
already follows from the place above $2$.
\end{proof}



\section{Exact three-adic depth of the anchor}
\label{sec:three-adic-anchor-depth}

The coefficient size of the physical anchor can be reduced by using its
exact local depth.  The following calculation concerns only the fixed
point and its integral multiples; it is independent of the prime-pattern
and global height hypotheses.

Put
\[
 K=\mathbb Q(\sqrt2),\qquad D=11-6\sqrt2,\qquad
 E:y^2=x^3-Dx,
\]
and
\[
 P=\left(\frac{57}{8}-\frac72\sqrt2,
          -\frac{165}{8}+\frac{417}{32}\sqrt2\right).
\]
The prime three is inert in $K$.  Let $v$ be the valuation on
$K_3=\mathbb Q_3(\sqrt2)$ normalized by $v(3)=1$.
As usual, the valuation of zero is $+\infty$.

\begin{proposition}\label{anchor3:exact-depth}
For every nonzero integer $c$, the point $[c]P$ is finite and has
nonzero abscissa.  Write
\[
             f([c]P)=\frac1{x([c]P)}=a_c+b_c\sqrt2.
\]
Then
\[
 \min\{v_3(a_c),v_3(b_c)\}=
 \begin{cases}
 0,&c\text{ odd},\\
 2+2v_3(c),&c\text{ even}.
 \end{cases}
\]
In the even case the formal parameter $z=-x/y$ satisfies
\[
                  v(z([c]P))=1+v_3(c).
\]
\end{proposition}
\begin{proof}
The displayed equation has good reduction at three: its discriminant is
$64D^3$, which is a unit.  Direct reduction gives
\[
             \overline P=(\sqrt2,0)\in E(\mathbb F_9).
\]
This is a nonidentity point of order two, with unit abscissa.  Hence
every odd multiple of $P$ is finite with unit abscissa.

The ordinate of $P$ has valuation one, since
\[
 y(P)=\frac{3(-220+139\sqrt2)}{32}
\]
and its parenthesized factor has nonzero residue in $\mathbb F_9$.
The numerator of the tangent slope is a unit: modulo three,
$3x(P)^2-D=-D\ne0$.  Thus
\[
 \lambda=\frac{3x(P)^2-D}{2y(P)}
 \quad\hbox{satisfies}\quad v(\lambda)=-1.
\]
The duplication formula now gives
\[
 \begin{split}
 v(x([2]P))&=v(\lambda^2-2x(P))=-2,\\
 v(y([2]P))&=v\bigl(\lambda(x(P)-x([2]P))-y(P)\bigr)=-3.
 \end{split}
\]
There is no cancellation: the summands in the first difference have
valuations $-2,0$, and those in the second have valuations $-3,1$.
Consequently $v(z([2]P))=1$.

For completeness, the needed formal-group valuation is elementary here.
The formal group of the integral Weierstrass equation has integral
coefficients, and negation sends $z$ to $-z$.  Its multiplication series
therefore have the form
\[
                  [m](Z)=mZ+Z^3g_m(Z^2),
                  \qquad g_m(T)\in\mathcal O_{K_3}[[T]].
\]
If $v(z)\ge1$ and $3\nmid m$, the linear term has strictly smaller
valuation than every higher term, so $v([m]z)=v(z)$.
For $m=3$, the linear term has valuation $v(z)+1$, whereas all higher
terms have valuation at least $3v(z)>v(z)+1$.  Thus
\[
                    v([3]z)=v(z)+1.
\]
Write $c=2\cdot3^ru$, where $3\nmid u$.  Repeated application to
$z([2]P)$ yields $v(z([c]P))=1+r$.  In particular, this parameter
is nonzero, so $[c]P$ is not the identity.  On the formal group,
$x=z^{-2}$ times an integral power series with constant term one.
It follows that
\[
                    v(f([c]P))=2+2r.
\]

Finally, $1,\sqrt2$ form an integral basis for the unramified quadratic
extension $K_3/\mathbb Q_3$, with linearly independent residues over
$\mathbb F_3$.  For any rational $a,b$, not both zero, this gives the
exact equality
\[
                  v(a+b\sqrt2)=\min\{v_3(a),v_3(b)\}.
\]
Applying it to the reciprocal abscissas proves both cases.
\end{proof}

\begin{corollary}\label{anchor3:integer-cutoff}
Let $N\ge1$ and $c\ne0$.  Then
\[
 \frac{a_c}{3^N},\frac{b_c}{3^N}\in\mathbb Z_3
 \quad\Longleftrightarrow\quad
 c\text{ is even and }v_3(c)\ge\left\lceil\frac N2\right\rceil-1.
\]
For a prescribed integer $r\ge0$, the smallest positive $N$ forcing
$v_3(c)\ge r$ is $N=2r+1$.  At this value the anchor implies
\[
                             |c|\ge2\cdot3^r.
\]
\end{corollary}
\begin{proof}
The displayed divisibility condition says that the minimum valuation in
Proposition~\ref{anchor3:exact-depth} is at least $N$.  Since $N>0$,
this excludes the odd case.  In the even case it is exactly
$2+2v_3(c)\ge N$.  The stated rounding follows.  When $r\ge1$,
the value $N=2r$ still admits $c=2\cdot3^{r-1}$, so $2r+1$ is
minimal; for $r=0$, it is the smallest positive value by definition.
The final bound follows from the simultaneous divisibility by two and
$3^r$.
\end{proof}

\section{A scalar normalization without squareclass corrections}
\label{sec:scalar-correction-free}

Put
\[
                         \beta(X)=216X^4+450459
\]
and retain the five-contact polynomial
\[
 r(T)=T(T-1)(T-2401)(32T+343)(2T-1).
\]
Define
\[
 r_X=r(\beta(X)),\qquad
 \delta_X=1+(17r_X)^2,\qquad
 R_X=\frac{102^2r_X}{\delta_X^2}.
\]
The scaling factor is an integer.  In particular, this normalization
introduces no fixed denominator prime.

\begin{lemma}\label{sn:integral-input}
For $X\in\mathbb Q$ one has
\[
                 \beta(X)\in\mathbb Z
                 \quad\Longleftrightarrow\quad X\in\mathbb Z.
\]
For every real $X$, $\beta(X)\ge450459>2401$, so $r_X>0$
and $\delta_X>0$.  If $p>76176$ and $v_p(\beta(X))<0$,
then
\[
 v_p(\beta(X))=4v_p(X),\qquad
 v_p(r_X)=5v_p(\beta(X)),\qquad
 v_p(R_X)=-3v_p(r_X).
\]
In particular, all three valuations are even.
\end{lemma}
\begin{proof}
If $X$ is an integer, so is $\beta(X)$.  Conversely, write
$X=a/b$ in lowest terms with $b>0$.  Integrality of $\beta(X)$
gives $b^4\mid216=2^3\cdot3^3$.  Neither prime exponent
admits a nontrivial fourth power, so $b=1$.

The real inequalities follow from the fourth power and the
five positive factors of $r(T)$ when $T>2401$.  At a prime
in the stated range, a negative valuation of $\beta$ forces
$v_p(X)<0$.  The unique dominant terms in $\beta$ and
$r(\beta)$ are respectively $216X^4$ and $64\beta^5$.
This gives the first two valuation formulas.  In this case
$v_p(r_X)<0$, so the dominant term of $\delta_X$ is
$17^2r_X^2$, and $v_p(\delta_X)=2v_p(r_X)$.  The final
formula follows.  The displayed multiple of four is even.
\end{proof}

\begin{lemma}\label{sn:local-squares}
If $X\in\mathbb Z$, then $r_X$ is a nonzero square in
$\mathbb Q_2$, $\mathbb Q_3$, and $\mathbb Q_{17}$.
More precisely,
\[
 v_2(r_X)=2,\quad r_X/4\equiv1\pmod8,
 \qquad
 v_3(r_X)=2,\quad r_X/9\equiv1\pmod3,
\]
and at seventeen either $v_{17}(r_X)=0$ with square-unit
residue, or
\[
                    v_{17}(r_X)=2,\qquad
                    r_X/17^2\equiv9\pmod{17}.
\]
\end{lemma}
\begin{proof}
Since $216$ is divisible by eight and by twenty-seven,
\[
                    \beta\equiv3\pmod8,
                    \qquad\beta\equiv18\pmod{27}.
\]
For the two even factors of $r(\beta)$ one has
\[
 v_2(\beta-1)=v_2(\beta-2401)=1,
 \qquad
 \frac{\beta-1}{2}-\frac{\beta-2401}{2}=1200\equiv0\pmod8.
\]
Their odd parts therefore have square product congruent to one
modulo eight.  The remaining odd factors have product
$3\cdot7\cdot5\equiv1\pmod8$.  This proves the assertions
at two.  At three, only the factor $\beta$ is divisible by
three; it has valuation two.  Reduction after division by nine
gives
\[
 \frac{r(\beta)}9\equiv
       2\cdot2\cdot2\cdot1\cdot2\equiv1\pmod3.
\]

At seventeen the fourth powers are exactly $0,1,4,13,16$.
Since $216\equiv12$ and $450459\equiv10\pmod{17}$, all
possibilities are given by the table
\[
\begin{array}{c|rrrrr}
 X^4&0&1&4&13&16\\\hline
 \beta&10&5&7&13&15\\
 r(\beta)&0&15&2&13&2.
\end{array}
\]
The nonzero residues are squares: $13=8^2$, $15=7^2$, and
$2=6^2$ modulo seventeen.

In the remaining case $17\mid X$.  The exact identity
\[
 32\cdot450459+343=17^2\cdot49879,
 \qquad 49879\equiv1\pmod{17},
\]
and $v_{17}(X^4)\ge4$ show that the factor
$32\beta+343=6912X^4+17^2\cdot49879$ has valuation exactly
two and unit residue one.  The other four factors are units,
with product
$10\cdot9\cdot6\cdot2\equiv9\pmod{17}$.
This proves the last assertion.  The usual local square criteria
at two and at odd primes now prove all three square statements.
\end{proof}

\begin{proposition}\label{sn:prepared-constant}
For every integer $X$, the number $R_X$ is prepared: it is
positive and less than one, its reduced denominator is a square,
it is a square in $\mathbb Q_2,\mathbb Q_3,\mathbb Q_{17}$,
and
\[
 v_2(R_X)=4,\qquad v_3(R_X)=4,\qquad
 v_{17}(R_X)\in\{2,4\}.
\]
No squareclass correction is required.
\end{proposition}
\begin{proof}
The contact value $r_X$ is a positive integer.  One has
$\gcd(r_X,\delta_X)=1$, and the preceding lemma shows that
$\delta_X$ is a unit at two and three; it is a unit at
seventeen by its definition.  Thus
$\gcd(102^2r_X,\delta_X^2)=1$, and the reduced denominator
is exactly $\delta_X^2$.

The quotient $R_X/r_X=(102/\delta_X)^2$ is a rational square.
Its local valuations add two at each of $2,3,17$, proving
both the asserted squareclasses and the exact depths.  Finally,
since $r_X\ge1$,
\[
 0<R_X=
 \frac{102^2r_X}{(1+17^2r_X^2)^2}
 <\frac{102^2}{17^4r_X^3}
 \le\frac{10404}{83521}<1.
\]
\end{proof}

\subsection{The norm tree and its domain}

Fix a positive rational number $\rho$ and six rational tests
$t_0,\ldots,t_5$, with their assigned positive integer orders.
Start the scalar norm tree with the distinguished node
$1+t_0^2/\rho$.  Its other initial nodes are the assigned
normalized tests or pair norms.  Every later operation is
\[
                  N_\rho(u,v)=u^2+\rho v^2
\]
or the unchanged carrying of an unpaired node.  This is the
constant-left tree used by the scalar certificate; denote its
root by $W_\rho$.

\begin{lemma}\label{sn:positive-root}
Where its rational expressions are defined, $W_\rho>0$.
When all six tests tend to zero, $W_\rho$ tends to one.
For fixed nonzero rational $\rho$, at any finite set of primes
the tests can be required sufficiently small that
$W_\rho$ is a unit at every prime in that set.
\end{lemma}
\begin{proof}
The distinguished initial node is positive.  At each norm
operation its positive descendant contributes a positive
square, with a positive coefficient if it is the second
argument.  Its path therefore remains positive up to the
root.  When all tests vanish, this path starts at one and
all other nodes vanish, so its root is one.  For fixed
$\rho$, the tree is a polynomial in the normalized tests
with fixed rational coefficients.  It is continuous at zero
over $\mathbb R$ and over every $\mathbb Q_p$.
Its value at zero is one.  The inverse image of $\mathbb Z_p^*$
is therefore a neighborhood of zero, proving the final claim.
\end{proof}

There are two valid implementations.  One may replace the contact
parameter throughout the tree by $\rho=R_X$ and put
$b=R_XW_{R_X}^2$.  Alternatively, one may retain the original
tree with $\rho=r_X$ and normalize only its final pole input:
\[
                       b=R_XW_{r_X}^2.
\]
The second choice retains the original test denominators.  It
also has the quadratic scaling equations
\[
        \delta_X=1+17^2r_X^2,\qquad
        \delta_X V=102W_{r_X},\qquad b=r_XV^2,
\]
where the last expression can use the existing square-and-product
gates.  The first two displayed equations are quadratic.
The soundness and completeness proof for this substitution is
given below.

For rational inputs, both implementations have $r_X,R_X,\delta_X>0$
and $W_\rho>0$.  Thus the scalar contact, normalizer, norm-root,
and pole-input nonvanishing conditions follow on the stated
rational-expression domain.  A coordinate used solely to certify
one of these nonzero conditions can be removed after its occurrence
condition is checked.  This observation does not remove unrelated
denominator guards in the elliptic-point formulas.

\subsection{Soundness and simultaneous completeness}

\begin{proposition}\label{sn:contact-preservation}
At every finite scalar contact prime $p>76176$, the replacement
has the same positive odd contact valuation as the original
certificate.  Consequently it preserves its norm-tree extraction
of the assigned contact orders, for either implementation above.
\end{proposition}
\begin{proof}
At such a contact, $\nu=v_p(r_X)>0$ is odd.  The normalizer
satisfies $\delta_X\equiv1\pmod p$, and $102$ is a unit.
Hence
\[
               v_p(R_X)=v_p(r_X)=\nu,
 \qquad R_X/r_X\in\mathbb Z_p^{*2}.
\]
For either $\rho=r_X$ or $R_X$, the two terms in
$u^2+\rho v^2$ have valuations of opposite parity when both
are nonzero.  There is therefore no cancellation: its valuation
is the minimum of $2v_p(u)$ and $\nu+2v_p(v)$.
Every step of the existing constant-left extraction uses exactly
this identity, the assigned powers of the contact parameter in
the leaves, and the odd-pole exclusion for $b$.  All these
valuations are unchanged.  Lemma~\ref{sn:integral-input}
also shows that a contact with infinity has even valuation
outside the fixed exceptional set, so it creates no omitted
odd contact.  The inherited height and comparison arguments
therefore give $\beta\in\mathbb Z$, and
Lemma~\ref{sn:integral-input} gives $X\in\mathbb Z$.
\end{proof}

\begin{proposition}\label{sn:simultaneous-completeness}
Under the same arithmetic inputs as the source construction,
the scalar block with the input map and normalization above
can be satisfied at every integer input, together with the
remaining point, height, comparison, and size conditions.
The finite squareclass correction parameter can be omitted.
\end{proposition}
\begin{proof}
Fix an integer input.  The numbers $r_X,R_X,\delta_X$ are
now fixed, before any contact primes are chosen.  Proposition~
\ref{sn:prepared-constant} supplies the sign, local squareclasses,
square denominator, and small positive real size formerly obtained
by a correction.  Include the finite support of the numerator
and denominator of $R_X$, and that of $\rho$, in the fixed
local approximation conditions for the common point $H$ and
the scalar point $S-P$.

The required approximation is the same one used in the source.
Choose $H$ sufficiently deep at these finite places and in a
small real neighborhood of the identity, excluding the identity
itself.  Shrink the real neighborhood so that its formal
parameter is bounded away from zero.  Choose $S-P$ still deeper
at every fixed finite place and sufficiently close to the real
identity.  The reciprocal-abscissa tests at these points then
tend to zero, and the quotient of their formal parameters
tends to zero with $H$ in that smaller neighborhood.
Fixed powers of $\rho$ in the denominators do not affect this
convergence.

One can impose these conditions with the required congruences
on the integer indices: the source first chooses the fixed
modulus and depths, then applies its prime-pattern and common-point
construction.  Increasing the modulus to include the finite
support of $\delta_X$ is allowed because this support was fixed
before that application.  It introduces no new family of
varying-prime conditions and no stronger prime-pattern theorem.
The real compact-group argument permits the same shrinking of
the neighborhoods and arbitrarily large scalar index.  All
other source index and comparison choices are unchanged.

By Lemma~\ref{sn:positive-root}, these choices make
$W_\rho$ a unit at every prime in the support of $R_X$ and
give $|W_\rho-1|<1/2$ at the real place.  Therefore
\[
                       0<b=R_XW_\rho^2<9/4<17.
\]
The factor $W_\rho^2$ preserves the local squareclasses and,
at $2,3,17$, preserves the depths of Proposition~
\ref{sn:prepared-constant}.  The reduced denominator of $b$
is a square: at a prime in the support of $R_X$, the root is
a unit, while at every other prime its valuation contribution
is even.  Thus $b$ is prepared.  The same inherited rational
pole-realization theorem supplies the final pole witnesses.
This proves completeness without a correction branch.
\end{proof}

The input polynomial has degree four and can be evaluated with
two quadratic gates while preserving the source expression
$\beta=2Y+3$:
\[
                    Z=X^2,\qquad Y=108Z^2+225228.
\]
The gain in any compiled polynomial also depends on its exact
occurrence and denominator inventory; the preceding propositions
establish the arithmetic substitution independently of that count.

\subsection{Number-field scope}

All displayed algebraic identities remain valid in characteristic
zero.  Positivity of the contact and norm root also holds in a
number field with a real embedding, after applying that embedding.
The local squareclass table, however, concerns rational integer
inputs and the three specified completions of $\mathbb Q$.
It does not establish the corresponding preparation assertions
for every algebraic integer in a larger field.

Even the integrality equivalence of Lemma~\ref{sn:integral-input}
does not automatically become an equivalence for rings of
integers.  Let $x_0$ be a real fourth root of $1/108$ and
$L=\mathbb Q(x_0)$.  Then
\[
                       \beta(x_0)=450461\in\mathbb Z.
\]
But $x_0$ is not an algebraic integer, since $x_0^4=1/108$.
Thus this elementary input map
alone does not transfer a definition to $\mathcal O_L$.
No assertion that the complete rational defining polynomial
defines every $\mathcal O_L$ is made or follows from this
normalization.


\section{Exact specifications and reproducibility}
The complete coefficient expression graphs and finite algebra relations
are supplied in the computation directory. Independent checkers verify
the rational source expressions, every monic relation, the guard
coverage, the finite alternatives, and the successive degree and
coefficient bounds. The producers reproduce the saved specifications
exactly. The indexed exclusion products and determinants are specified
by finite algorithms; their enormous complete expansions are not
needed for either definition or verification of the bounds.

The source arithmetic inputs are separate mathematical premises.
The reproduction checks certify the algebraic construction and the
stated quantitative implications, not those global arithmetic inputs.
\begin{thebibliography}{99}
\small
\bibitem{RepositoryV9}
E.~Glazer and contributors, \emph{Diophantine definition of the integers
in the rationals}, source repository, coefficient baseline
\texttt{14e8f601fe0d9173d4a5f084aedc3e07066a7688}, 7 November 2026.
The odd-product identity, direct targets and asymmetric-depth
construction are in \texttt{docs/quartic-523-540/addition.tex}; the
repository attributes the odd-product contribution to G.~Conant,
developed with ChatGPT.
\url{https://github.com/elliotglazer/diophantine-z-q/tree/14e8f601fe0d9173d4a5f084aedc3e07066a7688}.


\bibitem{Sun2026Rational86}
Z.-W.~Sun, \emph{86 unknowns for Hilbert's tenth problem over
the rationals}, arXiv:2610.20691v2, 19 October 2026,
especially Theorem 4.1 and Sections 10.3--10.6.
\url{https://arxiv.org/abs/2610.20691v2}.


\bibitem{Voight2021}
J.~Voight, \emph{Quaternion Algebras}, Graduate Texts in
Mathematics 288, Springer, 2021, especially Proposition 14.6.7.
\url{https://jvoight.github.io/quat.html}.


\bibitem{Silverman2009}
J.~H.~Silverman, \emph{The Arithmetic of Elliptic Curves},
second edition, Graduate Texts in Mathematics 106, Springer, 2009.
\url{https://doi.org/10.1007/978-0-387-09494-6}.


\end{thebibliography}
\end{document}
